% s-338-p338-m003; original PDF 338, 339
\Aphorism{s-338-p338-m003}{388}
Bin Ich es nicht, der die tausend Nadeln der Pinie zusammenhält und aus Millionen Wellen ein Meer glättet? Continuität des Bewusstseins projicirt sich als Consistenz der Aussenwelt; in dem regellosen Wirbel der Molecüle erliest sich das Subject sein Object, das Chaos wird durchgesiebt zum Kosmos. Die festen Linien der Umwelt sind die Fäden, an denen sich das Bewusstsein entlang tastet; die Gesetzlichkeit, die ich um mich vorfinde, ist Widerschein der Bedingung, unter der ich mich vorfinde. In dem einen Chaos sind viele \Greek{κόσμοι} möglich, wie im Raum viele Geraden und Kreise; jeder Ausschnitt aus dem Chaos, jede Beschränkung der transcendenten Willkür kann einer Abfolge von Bewusstseinsvorgängen entsprechen. Die Richtung einer Beschränkung festhalten, ist eine weitere Beschränkung, also gesteigertes Bewusstsein; dass Ich mir die Mühe nehme, die Pinie dauern zu sehen, trägt zur Befestigung meines Ich bei. Und je mehr Gesetzlichkeit, desto mehr Ich; je strengere Auslese unter dem Möglichen, um so reichere Wirklichkeit; je enger das Sieb, um so feiner das Korn.
\par

